\chapter{Glosario}
\label{app:glossary}
% ==============================================================

\begin{description}[leftmargin=3cm, style=nextline]
  \item[EPSG:25830] ETRS89 / UTM zona 30N. Sistema de referencia cartográfico oficial para la Península Ibérica.
  \item[GCP] \textit{Ground Control Point}. Punto de control terrestre con coordenadas conocidas.
  \item[Homografía] Transformación proyectiva planar entre dos planos (imagen y terreno).
  \item[IEL] Instituto de Ecología Litoral. Suministra los GCPs GNSS del proyecto.
  \item[RANSAC] \textit{Random Sample Consensus}. Estimador robusto de parámetros geométricos.
  \item[ROI] \textit{Region of Interest}. Zona de la imagen relevante para el análisis.
  \item[SAM] \textit{Segment Anything Model}. Modelo de segmentación de Meta AI.
  \item[SIFT] \textit{Scale-Invariant Feature Transform}. Detector de puntos de interés invariante a escala y rotación.
  \item[RMSE] \textit{Root Mean Square Error}. Error cuadrático medio.
  \item[MAE] \textit{Mean Absolute Error}. Error medio absoluto.
  \item[Farneback] Algoritmo de flujo óptico denso para estimar campo de desplazamiento entre dos imágenes.
  \item[Inliers] Correspondencias que se ajustan al modelo estimado por RANSAC.
  \item[Two-phase commit] Patrón de escritura en dos fases: staging temporal → confirmación permanente.
  \item[Drift] Desplazamiento residual de píxeles tras la corrección de homografía.
  \item[Auto Mode] Módulo que automatiza, para un rango de fechas elegido, la secuencia
        descarga → alineación → análisis sobre una cámara ya calibrada
        (capítulo~\ref{ch:automode}).
  \item[Aviso de geometría inestable] (\texttt{geometry\_warning}) Mensaje que advierte de que
        RANSAC no encontró un subconjunto mutuamente consistente de varillas suficientemente
        grande al calcular la homografía — el RMSE resultante puede no reflejar la calidad real
        de la calibración (capítulo~\ref{ch:calibration}).
  \item[Varillas usadas] Número de GCPs que entran en el ajuste de la homografía (confirmadas y
        no excluidas) — \textbf{no} es lo mismo que \textit{inliers}, que es cuántas de esas
        RANSAC clasifica como mutuamente consistentes (capítulo~\ref{ch:calibration}).
\end{description}

% ==============================================================
